\documentclass[conference]{IEEEtran}
\IEEEoverridecommandlockouts

\usepackage{cite}
\usepackage{amsmath,amssymb,amsfonts}
\usepackage{algorithmic}
\usepackage{graphicx}
\usepackage{textcomp}
\usepackage{xcolor}
\usepackage{booktabs}
\usepackage{hyperref}
\usepackage{float}
\usepackage{chngcntr}

\renewcommand{\thesection}{\arabic{section}}
\counterwithin{figure}{section}
\def\BibTeX{{\rm B\kern-.05em{\sc i\kern-.025em b}\kern-.08em
    T\kern-.1667em\lower.7ex\hbox{E}\kern-.125emX}}
\renewcommand{\thesubsection}{\thesection.\arabic{subsection}}
\begin{document}

\title{Agro-Veritas: Hallucination-Aware Coffee Disease Diagnosis using Retrieval-Augmented Generation}

\author{
\small
\begin{minipage}[t]{0.23\textwidth}\centering
\textbf{Vaidurya Samaga}\\
\textit{Department of Computer Science}\\
\textit{PES University}\\
Bengaluru, India
\end{minipage}\hspace{0.01\textwidth}
\begin{minipage}[t]{0.23\textwidth}\centering
\textbf{Vishwas H Gowda}\\
\textit{Department of Computer Science}\\
\textit{PES University}\\
Bengaluru, India
\end{minipage}\hspace{0.01\textwidth}
\begin{minipage}[t]{0.23\textwidth}\centering
\textbf{Kusumita}\\
\textit{Department of Computer Science}\\
\textit{PES University}\\
Bengaluru, India
\end{minipage}\hspace{0.01\textwidth}
\begin{minipage}[t]{0.23\textwidth}\centering
\textbf{Dr.Prajwala TR}\\
\textit{Associate Professor}\\
\textit{Department of Computer Science and Engineering}\\
\textit{PES University}\\
Bengaluru, India
\end{minipage}
}

\maketitle

\begin{abstract}
Retrieval-Augmented Generation (RAG) is increasingly used in the agricultural field such as disease diagnosis, but it still faces some important issues. Even if they may not accurately fit the specified symptoms, some diseases show up more frequently in the retrieved results. Additionally, the system may produce inaccurate or results (hallucinations) if farmers supply inadequate or ambiguous data. In this work, a balanced hybrid RAG-based system is proposed for coffee disease diagnosis to address these challenges. The system combines information stored in JSON format with unstructured knowledge from agricultural PDF documents, ensuring better coverage and balance across the diseases. The system also includes multi-turn interaction, which enables the system to ask simple clarification questions to the user. In addition, multiple responses are generated and compared with each other to measure consistency using a hallucination checking module. 

The system was evaluated on a small labeled dataset and it achieved a Top-1 accuracy of 50\% in the single-turn setup and improved to 66.67\% when multi-turn interaction was used. The Top-3 accuracy was 83.33\% which was used as a proxy measure. The average consistency score was 88.89\% while the hallucination rate was 25\%. Keyword-based approach was used to measure Precision@5 which came to 48.19\% and the mean CRAG relevance score was 0.6588. Overall, the proposed approach improves the reliability of RAG-based diagnosis systems for real-world agricultural use.

\end{abstract}

\begin{IEEEkeywords}
retrieval-augmented generation, hallucination detection, multi-turn dialogue, disease diagnosis, hybrid retrieval, ambiguity detection, corrective RAG, agricultural AI
\end{IEEEkeywords}

\section{INTRODUCTION}

Diagnosis of coffee plant diseases must be accurate for ensuring good yield and quality of the crops especially in regions like Karnataka, where coffee cultivation plays a significant role in the economy. In practice, farmers usually describe the symptoms of the plant informally without detailed characterisation. This may lead to an inaccurate diagnosis. Hence, effective diagnosis depends on both the quality and specificity of the input provided. A critical observation made is that a single symptom rarely corresponds to a specific disease. For example, yellow spots on leaves are a common symptom in multiple coffee diseases, making diagnosis ambiguous without additional details. Relying only on the incomplete information given as input by the farmer leads to incorrect diagnosis and wrong treatment recommendations. This makes it necessary for the system to recognise ambiguity, seek clarification, and provide evidence-based recommendations before concluding about disease diagnosis.

Recent advancements in digital agriculture have allowed the use of advisory tools assisted with large language models and Retrieval-Augmented Generation(RAG) methods. In such systems, responses are verified using grounded documents that are domain-specific. However, most of the existing systems operate in a single-turn setting, assuming a single user query is sufficient. In practice, farmer queries are often ambiguous or incomplete. When such queries are processed by even RAG-based systems, they may generate responses based on partially relevant information, producing confident but false recommendations---the hallucination problem.

Furthermore, RAG systems exhibit systematic failure modes in agricultural diagnosis. Single-disease bias occurs when certain diseases dominate retrieval regardless of actual relevance to symptoms. This bias emerges because dense retrievers learn to associate common symptom keywords with retrievable documents, causing the most-documented disease to accumulate retrieval probability. When Coffee Leaf Rust is well-represented in training documents, systems tend to retrieve it for nearly all yellow-leaf symptoms, suppressing evidence for less-documented but equally plausible diseases.

To address these challenges, we propose a hallucination-aware, multi-turn coffee disease diagnosis system that leverages RAG with explicit ambiguity detection, evidence-grounded clarification, and hallucination verification. The system retrieves information from verified agricultural documents and ensures retrieved information is used both in intermediate and final outputs. The key innovations are ambiguity detection in farmer queries to prevent premature diagnosis, hybrid knowledge retrieval combining structured JSON disease knowledge (60\%) with unstructured PDF evidence (40\%) to prevent single-disease dominance, evidence-grounded clarification question generation to gather missing symptom information, multi-turn interactive diagnosis that improves accuracy through iterative refinement, reliable, context-aware treatment recommendations grounded in retrieved evidence, explicit hallucination detection comparing consistency across multiple independent generations.

This system demonstrates that balanced dataset design, complemented by multi-turn interaction and explicit verification, enables trustworthy agricultural diagnosis while maintaining computational efficiency through local inference. Our approach provides a replicable framework for any domain-specific diagnosis system requiring both accuracy and verifiable trustworthiness.

The contributions of this work are the identification and mitigation of single-disease bias through balanced hybrid dataset architecture, an ambiguity-aware multi-turn diagnostic pipeline with evidence-grounded clarification questions, a hallucination detection mechanism using consistency verification across multiple generations, a comprehensive methodology for agricultural disease diagnosis requiring document-grounded evidence, and empirical evaluation demonstrating elimination of retrieval bias while maintaining diagnostic accuracy.

\section{RELATED WORK}

The field of Retrieval-Augmented Generation(RAG) has grown a lot as a way to improve the answers given by large language models(LLMs). Instead of depending only on what the model already knows, RAG connects the model to external data sources. This helps reduce wrong answers and makes the output more accurate. Recent work mainly focuses on using RAG in specific domains, handling errors during retrieval, and checking whether the generated answers are correct.

In this area, Samuel et al. [1] developed AgroLLM to connect agricultural research with real farming needs. Their system uses textbooks and research papers stored in FAISS database to quickly find useful information. Their results showed that using RAG improves performance, and ChatGPT-4o mini achieved around 93\% accuracy for agricultural queries. Similarly, Singh et al. [2] developed Farmer Chat,a system designed to help small farmers. It uses a multi-step RAG process and can handle different types of data like text, tables, and videos. This system has bee used in different countries and has helped more than 15,000 farmers. While AgroLLM focuses more on model performance, Farmer.Chat focuses more on user needs like trust, easy understanding , and accessibility.

Some researchers have worked on improving RAG when the first retrieval is not correct. The CRAG framework [3] uses a simple evaluator to check whether the retrieved documents are useful or not. If the results are not good, the system searches the web to find better information. This method improved performance by 19\% on the popQA dataset .In another work, Kebir et al. [4] proposed the RAC framework to handle unclear the final answer is more accurate. CRAG mainly focuses on fixing wrong retrieval, while RAC focuses on better understanding the user query.

Another challenge is detecting wrong or made-up answers. Manakul et al. [5] proposed SlefCheckChatGPT, which checks if the model gives consistent answers. It generates multiple responses and compares them. If the answers are different, it means the model may be giving incorrect information. This method works without needing external database.

These studies show the RAG systems are becoming more accurate and useful.However, some methodsare slow so notfully consider real-world user needs. To solve this, our work combines the strengths of CRAG with the user-focused approach of Farmer.Chat. This help in building a system that is both accurate and easy to use in real situations.

\section{SYSTEM ARCHITECTURE}

\subsection{Overview}

Our system is designed to help farmers identify coffee plant diseases in a simple and reliable way.It uses multiple components that work together step by step. The system follows a multi-turn approach, meaning it can ask follow-up questions to improve the final result. The architecture diagram of the system is as shown in Fg. 3.1.

\begin{figure}[H]
    \centering
    \includegraphics[width=\columnwidth]{image.png}
    \caption{Architecture Diagram}
    \label{fig:architecture}
\end{figure}

The first component is the \textbf{PDF Loader},which collects information from three trusted PDF documents.These documents are split into smaller parts so that the system can process them easily.

Next, the \textbf{Vector Store(FAISS)} is used to store and search this information efficiently.It helps in quickly finding relevant data based on the user’s input.

We also use a \textbf{JSON Retreiver}, which contains structured data about different coffee diseases, including symptoms, treatment, and prevention. This helps in giving more accurate and clear results.

The \textbf{Hybrid Retriever} combines both PDF and JSON data.Instead of relying on only one source, it uses both(60\% JSON and 40\% PDF) to improve accuracy and reduce bias.

To combine results from both sources, we use the following scoring formula:

\textit{\textbf{score(d) = 0.6⋅sJSON(d) + 0.4⋅sPDF(d)}}

This means the system gives more importance to JSON data because it is structured, while PDF data provides additional supporting information.

The \textbf{Ambiguity Detector} checks whether the user’s input is incomplete. For example, if the user says “yellow leaves” the system identifies missing details like location or pattern.

The \textbf{Retrieval Evaluator(CRAG)} removes irrelevant or low quality results so that only useful information is used.

The \textbf{Clarification Generator} asks simple follow up questions to get missing details from the user.

The State Manager keeps track of the conversation, including previous responses and confidence level.

The Diagnosis Generator produces the final output, including the disease name, confidence score, and suggestions for treatment and prevention.

Finally, the Hallucination Checker verifies the result by generating multiple responses and comparing them.

\subsection{Hybrid Dataset Architecture}

The system uses a balanced dataset to avoid bias.

The JSON Knowledge base contains multiple coffee diseases from different categories, such as fungal, pest, and nutritional. Each disease is equally represented to ensure fairness.

The PDF documents provide detailed explanations and are stored in FAISS for fast retrieval.

By combining both sources, the system ensures that the output is both accurate and well-supported.

\subsection{Multi-turn Diagnostic Pipeline}

The system works in multiple steps.

First, the user gives symptoms. The system checks if the input is clear or if more information is needed. If details are missing, it asks follow-up questions.

Then, it retrieves relevant data using the hybrid retreiver and filters it using CRAG. This process continues for a few steps(maximum 3 turns) until enough information is collected.

Finally, the system generates the diagnosis and verifies it before showing the result.

\subsection{Stop Conditions}

The system stops asking questions when:

\begin{itemize}
    \item Confidence level is a high score (above 80\%)
    \item Only one disease remains
    \item Maximum number of questions(3) is reached
\end{itemize}
This helps balance accuracy and user convenience.

\subsection{Hallucination Verification}

After generating the result, the system checks whether it is reliable. It creates multiple responses and compares them.

To measure this, we use the following formula:

\textit{Consistency(y\^)=M1∑i=1M1[claims are consistent]}

\textit{If most of the responses are similar, the result is considered reliable.If they are different the system ask for more information or mark the result as uncertain.}

\section{METHODOLOGY}
In this paper, we present a hallucination-aware, multi-turn coffee disease diagnosis system leveraging the power of RAG. This system consists of an end-to-end pipeline that addresses ambiguity detection, hybrid retrieval, clarification questions, and hallucination detection for diagnosing the coffee disease. Some of the important components of this system are listed below.

\subsection{Stage 1: Processing of Input and Ambiguity Detection}

The system first accepts a natural language input from the user, describing the symptoms in coffee plants. Since farmer inputs are usually incomplete and lack sufficient details, an ambiguity detection module is implemented as the first step. Symptom-related attributes such as colour, pattern, location, and spread are extracted from the user’s input. The system uses rule-based parsing with light-weight semantic understanding to identify the attributes present in the input and the missing ones. For example, if the input is “leaves are turning yellow”, the system detects that the colour attribute is present and other attributes required for diagnosis, such as location(upper/lower leaves), pattern(uniform or spotted) and spread(localised or widespread) are missing. Each missing attribute is assigned a priority score, which is passed to the later stages for retrieval and clarification.

This step prevents the system from performing a diagnosis on insufficient or ambiguous input, thus reducing the chances of hallucinated outputs.

\subsection{Stage 2: Hybrid Knowledge Retrieval (RAG)}

To ensure evidence-grounded response generation, the system uses a hybrid Retrieval Augmented Generation (RAG) approach that combines structured and unstructured knowledge sources.

1) Structured Retrieval(JSON-based)

A curated dataset of coffee diseases is maintained in JSON format. Each disease entry includes the symptoms, causes, treatment methods, and prevention strategies. During retrieval, the system matches extracted symptom attributes with these structured entries.

2) Unstructured Retrieval(PDF-based)

Domain-specific agricultural documents are processed and indexed using FAISS, a vector database. Text chunks are embedded using a sentence piece transformer(all-MiniLM-L6-v2) to enable semantic similarity search.

3) Query Expansion 

Instead of having a single query, the system constructs multiple queries based on the original user input, extracted symptom keywords and the user responses to previous queries.

A comprehensive evidence pool is formed by combining structured and unstructured retrieval, which is used for evaluation.

\subsection{Stage 3: Retrieval Evaluation (CRAG)}

The Corrective Retrieval-Augmented generation(CRAG) mechanism is applied to refine the results, since the retrieval may contain irrelevant information. A relevant score is assigned to each retrieved document, which is computed using semantic similarity with the query and matches with the extracted attributes. Documents with scores lower than the threshold are discarded. Deduplication is used to eliminate redundant information and source diversity prevents over reliance on a single document. 

This stage ensures high-quality, relevant and diverse retrieval from evidence-grounded sources which directly improves the reliablility of question generation and diagnosis.

\subsection{Stage 4: Clarification Question Generation (RAC)}

Based on the missing attributes identified in Stage 1 and validated evidence from Stage 3, the system generates follow-up questions using an LLM.

The module for generating questions is built with the following constraints:-

\begin{enumerate}
    \item Single-focus questions: each one focuses on a single missing attribute.
    \item Evidence grounding: recovered context is used to formulate queries.
    \item Easy-to-understand language without scientific jargon to suit farmers
\end{enumerate}
For example, if location is a missing attribute, the system generates:

“Are the yellow leaves present on the lower part of the plant or the upper part?”

A significant advancement above conventional chatbot techniques is the use of evidence-based questioning, which ensures that the enquiries are context-aware and relevant to the disease rather than general.

\subsection{Stage 5: Multi-Turn Interaction}

The system supports interactive interaction where additional clarification questions are asked in order to collect more details for proper diagnosis.

A dedicated state manager module maintains:

\begin{itemize}
    \item previous user responses
    \item retrieved context across turns
    \item list of possible diseases
    \item confidence score for diagnosis
\end{itemize}
After every reponse from the user:

\begin{enumerate}
    \item The query is updated
    \item Retrieval and evaluation are repeated
    \item Candidate diseases are narrowed down
\end{enumerate}
The interaction terminates when one of the following conditions is met:-

\begin{itemize}
    \item Confidence score exceeds threshold
    \item Only one disease remains as the most possible candidate
    \item The maximum number of turns is reached
\end{itemize}
The multi-turn refinement allows the system to gradually converge on a correct diagnosis, which greatly decreases ambiguity.

\subsection{Stage 6: Disease Diagnosis Generation}

Once sufficient information is collected, the system returns the final diagnosis using an LLM.

The model is depends on:

\begin{itemize}
    \item filtered evidence from JSON and PDF sources
    \item the whole collection of user inputs
    \item interaction history throughout all turns
\end{itemize}
This approach guarantees that the output is based on retrieved knowledge, in contrast to direct LLM-based diagnosis.

The generated response consists of:

\begin{itemize}
    \item predicted disease name
    \item explanation based on symptoms
    \item recommended steps for treatment
    \item preventive measures
\end{itemize}
This phase, which combines retrieval-based evidence with generating capabilities, forms the system's core reasoning component.

\subsection{Stage 7: Hallucination Detection using Self-Consistency (SelfCheckGPT)}

To address the challenge of hallucinations, a self-consistency-based verification mechanism is applied. Instead of a single response, the system generates multiple diagnoses using the same input context. The outputs are compared across disease diagnosis, treatment recommendations and symptom interpretation.
Based on the level of agreement between the outputs, a confidence score is computed. If the consistency is high, output is considered to be reliable, else it indicates possible hallucination and the system may flag the result. 
This stage ensures robustness and reliability by serving as a safety layer.

\section{EXPERIMENTAL SETUP}

\subsection{Dataset and Baselines}

\textbf{Dataset}: 15 balanced coffee diseases from the structured JSON knowledge base, supplemented by three authoritative PDFs. The corpus ensures 6-7\% representation per disease, preventing dominance.

\textbf{Baseline Systems}:
\begin{itemize}
    \item \textbf{Standard RAG}: Single retrieval (top-5 documents), generate diagnosis, no clarification, no verification. Represents minimal viable RAG.
    \item \textbf{Always-Recursive RAG}: Three fixed retrieval-generation loops unconditionally, no verification, no clarification. Represents brute-force iteration.
    \item \textbf{Balanced Hybrid RAG} (Proposed): Multi-turn pipeline with hybrid retrieval, ambiguity detection, CRAG filtering, clarification, and verification.
\end{itemize}

All baselines share identical embedding models, vector stores, and generation components. Only orchestration differs.

\subsection{Evaluation Metrics}

\begin{table}[H]
\centering
\textbf{Table 6.1: Evaluation Metrics}\\[2pt]
\begin{tabular}{|p{2.5cm}|p{4.6cm}|}
\hline
\textbf{Metric} & \textbf{Definition} \\
\hline
Disease Coverage & \% representation of each disease in diagnoses (target: uniform $\approx$6.7\% per disease) \\
\hline
Hallucination Rate & \% of generated claims unsupported by retrieved context \\
\hline
Verification Accuracy & \% of diagnoses with consistency score $\geq 70\%$ \\
\hline
Avg. Turns to Diagnosis & Mean number of clarification rounds before termination \\
\hline
Claim Verification Rate & \% of generated claims verified as supported by context \\
\hline
User Satisfaction & Subjective rating of diagnosis clarity and confidence (1-5 scale) \\
\hline
\end{tabular}
\par\medskip
\noindent The performance of the proposed system is evaluated using the metrics in Table 6.1. It comprises metrics for diagnosis accuracy, hallucination detection, verification reliability, user satisfaction and interaction efficiency. These metrics are intended to measure the correctness of predictions as well as the reliability of the produced results.
\label{tab1}
\end{table}

\subsection{Hyperparameters}

\begin{table}[htbp]
\centering
\textbf{Table 6.2: Implementation Hyperparameters}\\[2pt]
\begin{tabular}{|p{3.8cm}|p{2.9cm}|}
\hline
\textbf{Component} & \textbf{Value} \\
\hline
Embedding Model & all-MiniLM-L6-v2 (384-dim) \\
\hline
PDF Chunk Size & 500 tokens \\
\hline
Chunk Overlap & 50 tokens \\
\hline
Hybrid Retrieval Ratio & 60\% JSON, 40\% PDF \\
\hline
CRAG Relevance Threshold & 0.3 \\
\hline
Max Turns ($T_{\max}$) & 3 \\
\hline
Confidence Threshold & 0.8 \\
\hline
Verification Consistency Threshold & 0.7 \\
\hline
Hallucination Check Generations & 3 \\
\hline
LLM Backend & Ollama (phi3) or Claude \\
\hline
\end{tabular}
\par\medskip
\noindent The system's key hyperparameters and design decisions are compiled in this Table 6.2. It contains information about the LLM backend, chunk size, embedding model, retrieval configuration, thresholds for CRAG filtering and verification levels. These parameters were selected to balance performance, efficiency, and reliability.
\label{tab2}
\end{table}

\section{RESULTS}

\subsection{Disease Coverage Analysis}

\textit{Hybrid balance testing on 8 diverse queries shows no single-disease dominance. Coffee Leaf Rust appeared in 2/8 outputs (25.0\%), while five other diseases were also predicted.}

\begin{table}[htbp]
\centering
\textbf{Table 6.3: Hybrid Disease Distribution}\\[2pt]
\begin{tabular}{|l|c|}
\hline
\textbf{Predicted Disease} & \textbf{Count (\%)} \\
\hline
Nitrogen Deficiency & 2 (25.0\%) \\
\hline
Coffee Leaf Rust & 2 (25.0\%) \\
\hline
Anthracnose & 1 (12.5\%) \\
\hline
Potassium Deficiency & 1 (12.5\%) \\
\hline
Brown Eye Spot & 1 (12.5\%) \\
\hline
Coffee Berry Borer & 1 (12.5\%) \\
\hline
\end{tabular}
\par\medskip
\noindent To determine whether the system is biased toward a single disease, Table 6.3 displays the distribution of predicted diseases over 8 test queries. The results show that the predictions are not biased, with no single disease dominating the output, demonstrating the effectiveness of hybrid RAG approach.
\label{tab3}
\end{table}

\subsection{Hallucination Reduction}

\textit{On the current labeled evaluation run, hallucination rate is 25.00\% with an average consistency score of 88.89\%.}

\begin{table}[htbp]
\centering
\textbf{Table 6.4: Evaluation Summary}\\[2pt]
\begin{tabular}{|p{3.8cm}|c|}
\hline
\textbf{Metric} & \textbf{Value} \\
\hline
Top-1 Accuracy (single-turn baseline) & 50.00\% \\
\hline
Top-1 Accuracy (multi-turn) & 66.67\% \\
\hline
Top-3 Accuracy (multi-turn, proxy) & 83.33\% \\
\hline
Avg Consistency Score & 88.89\% \\
\hline
Hallucination Rate & 25.00\% \\
\hline
Avg Precision@5 (keyword proxy) & 48.19\% \\
\hline
Avg CRAG relevance score & 0.6588 \\
\hline
Top-1 improvement (multi-turn - single-turn) & +16.67\% \\
\hline
\end{tabular}
\par\medskip
\noindent The system's overall performance across 12 test cases is shown in Table 6.4. It contains retrieval quality measurements, accuracy metrics for single-turn and multi-turn settings, consistency and hallucination ratings, and improvement resulting from multi-turn interaction. The results show that multi-turn clarification improves accuracy and helps in maintaining a good consistency in it's outputs.
\label{tab4}
\end{table}

\subsection{Diagnostic Efficiency}

\textit{Multi-turn interaction improved Top-1 accuracy by +16.67 percentage points over the single-turn baseline on the current 12-case benchmark.}

\subsection{Ablation Study}

\textit{A full ablation (component drop tests) is left as future work; this paper reports end-to-end system performance with hybrid retrieval, clarification, and verification enabled.}

\subsection{Qualitative Examples}

\textbf{Example 1 - Ambiguous Symptoms}:
\begin{itemize}
    \item \textit{User}:  ``Leaves turning yellow with brown spots''
    \item \textit{System}: ``On which parts of the plant are these yellow leaves appearing most? Lower leaves, upper leaves, or throughout?''
    \item \textit{User}: ``Mainly on lower leaves''
    \item \textit{System}: ``Are the brown spots forming patterns such as rings or spreading?''
    \item \textit{User}: ``They look like rings''
    \item \textit{System}: [Diagnosis] Coffee Leaf Rust, Confidence 87\%, Verified: 94\% consistency
\end{itemize}

The system correctly identified Coffee Leaf Rust by gathering information on colour, location and pattern.

\textbf{Example 2 - Ambiguous to Surprising Disease}:
\begin{itemize}
    \item \textit{User}: ``Small holes in leaves, some yellow edges''
    \item \textit{System}: (Clarification) ``Are the holes round or irregular?''
    \item \textit{User}: ``Perfectly round, very small''
    \item \textit{System}: [Diagnosis] Red Spider Mites, Confidence 76\%, Verified: 81\% consistency
\end{itemize}

This example shows how balanced retrieval and clarification might highlight uncommon pests rather than concentrating too much on one frequently occurring disease.

\section{DISCUSSION}

\subsection{Implications for Agricultural Diagnosis}

The results show that using a balanced dataset, along with multi-turn interaction and a verification step, helps in making the diagnosis more reliable. In the hybrid-balance test, Coffee Leaf Rust did not dominate the results and appeared in 25.0\% (2 out of 8) of the cases. The use of multi-turn interaction improved the Top-1 accuracy from 50.00\% to 66.67\%, which is an increase of 16.67 percentage points. The system also achieved an average consistency of 88.89\%, and the hallucination rate was 25.00\%.

These findings show that the SelfCheckGPT-inspired verification layer is an efficient way of lowering inaccurate outputs. Additionally, the clarification questions were crucial in addressing farmers' ambiguous inputs.

\subsection{Limitations}

\begin{itemize}
    \item \textbf{Limited Disease Set}: Evaluation on 15 diseases may not fully reflect performance on larger codebases (100+ diseases). Generalization to crops beyond coffee is untested.

    \item \textbf{Static PDF Corpus}: PDFs are not continuously updated with new disease variants or regional information. Dynamic knowledge integration would improve robustness.

    \item \textbf{Local LLM Performance}: Phi-3 and similar 7B-parameter models are less capable than GPT-4 or Claude-3. Results may not transfer to larger models with different hallucination profiles.

    \item \textbf{No Expert Validation}: Diagnoses are not compared against expert agronomists. Field validation is essential before deployment.

    \item \textbf{Limited User Study}: Evaluation is computational; subjective user satisfaction and trust metrics require larger user studies.
\end{itemize}

\subsection{Future Work}

\begin{itemize}
    \item Expand disease dataset to 100+ global coffee diseases and regional variants
    \item Integrate real-time farmer feedback to dynamically improve dataset balancing
    \item Validate diagnoses against ground-truth expert labels
    \item Conduct user studies with smallholder farmers in coffee-growing regions
    \item Extend framework to other crops (cocoa, cacao, tea)
    \item Deploy on edge devices with quantized models for offline access
    \item Integrate real-time image analysis for leaf symptom extraction
\end{itemize}

\section{CONCLUSION}

We have presented Balanced Hybrid RAG with Multi-turn Clarification and Hallucination Detection for coffee disease diagnosis. The system addresses three critical failures in standard RAG for agricultural diagnosis: single-disease bias through hybrid balancing, insufficient evidence through multi-turn clarification, and hallucinations through explicit verification.

Key results from current experiments are: (1) no single-disease dominance in hybrid testing, with Coffee Leaf Rust at 25.0\% (2/8); (2) Top-1 accuracy improvement from 50.00\% (single-turn) to 66.67\% (multi-turn); (3) Top-3 accuracy of 83.33\% (proxy); (4) average consistency of 88.89\% with a hallucination rate of 25.00\%.

Our work demonstrates that RAG-based diagnosis systems benefit from explicit bias mitigation at the dataset level, coupled with interactive clarification and verification. The framework is training-free and applicable to other agricultural, medical, and technical diagnostic domains where hallucination-free, trustworthy recommendations are paramount.

The combination of balanced retrieval, multi-turn interaction, and explicit verification provides a replicable recipe for deploying deep learning-based diagnosis systems responsibly in high-stakes domains where AI recommendations directly impact farmer livelihoods and food security.

\section*{ACKNOWLEDGMENT}

This work was supported by local computational resources and the Ollama open-source project for enabling efficient local LLM inference. The authors thank the open-source community for FAISS, Sentence-Transformers, and LangChain.

\begin{thebibliography}{00}

\bibitem{b1} D. J. S. Ravindran, I. Skarga-Bandurova, S. V., M. Awais, and S. M.,"AgroLLM: Connecting farmers and agricultural practices through large language models for enhanced knowledge transfer and practical application," 
AgriEngineering, vol. 8, no. 1, p. 38, 2026, doi: 10.3390/agriengineering8010038.

\bibitem{b2} N. Singh, J. Wang’ombe, N. Okanga, T. Zelenska, J. Repishti, J. G. K., S. Mishra, R. Manokaran, V. Singh, M. I. Rafiq, R. Gandhi, and A. Nambi,"Farmer.Chat: Scaling AI-Powered Agricultural Services for Smallholder Farmers," arXiv preprint arXiv:2409.08916, 2024, doi: 10.48550/arXiv.2409.08916.

\bibitem{b3} S.-Q. Yan, J.-C. Gu, Y. Zhu, and Z.-H. Ling,"Corrective Retrieval Augmented Generation," 
presented at ICLR (under review), 2024. [Online]. Available: https://openreview.net/forum?id=JnWJbrnaUE

\bibitem{b4}A. R.  Kebir, V. Guigue, L. S. Lhadj, and L. Soulier, "RAC: Retrieval-Augmented Clarification for Faithful Conversational Search," 
arXiv preprint arXiv:2601.11722, 2026, doi: 10.48550/arXiv.2601.11722.

\bibitem{b5} Y. El Kheir, A. Ali, and S. A. Chowdhury, 
"Automatic Pronunciation Assessment - A Review," 
in Findings of the Association for Computational Linguistics: EMNLP 2023, Singapore, Dec. 2023, pp. 8304–8324, doi: 10.18653/v1/2023.findings-emnlp.557.

\end{thebibliography}

\end{document}
